Une tige conductrice de longueur $L = \qty{\v{L}}{cm}$ glisse sans frottement à vitesse constante $v = \qty{\v{v}}{m/s}$ sur deux rails conducteurs parallèles. Les rails sont reliés à une extrémité par une résistance $R = \qty{\v{R}}{\ohm}$ ; la résistance de la tige et des rails est négligeable. Un champ magnétique uniforme de \qty{\v{B}}{T} est perpendiculaire au plan des rails.

\begin{enumerate}[label=\alph*)]
    \item Déterminez la force électromotrice induite.
    \item Déterminez le courant induit dans le circuit.
    \item Déterminez la force qu'il faut appliquer sur la tige pour maintenir sa vitesse constante.
    \item Montrez que la puissance fournie par cette force est égale à la puissance dissipée dans la résistance, et calculez-la.
\end{enumerate}
